% SPDX-License-Identifier: GPL-3.0-or-later
% Copyright (C) 2026 Alessandro Greco
%% Testo alternativo per l'accessibilita' (usato solo nell'EPUB).
\altfig{Nel caso storico del 2014, un HSDir ostile codifica il nome di un onion service nei tipi di cella RELAY e RELAY\_EARLY inviati verso il client. Un guard ostile dello stesso gruppo riconosce il segnale e associa il client alla richiesta di un descrittore, non a una visita dimostrata. Le quattro celle disegnate sono solo schematiche; i comandi sono visibili ai relay ma protetti da TLS sui collegamenti.}
% L'attacco del 2014: un segnale codificato nel tipo di cella, letto dal guard
\begin{tikzpicture}[font=\sffamily\scriptsize,
  ostile/.style={relay,fill=avviso!20,draw=avviso,align=center,font=\sffamily\scriptsize}]
  \node[client] (c) at (0,0) {client};
  \node[ostile] (g) at (3,0) {guard\\(ostile)};
  \node[relay] (m) at (6,0) {interm.};
  \node[ostile] (h) at (9,0) {HSDir\\(ostile)};
  \draw[<-,thick] (c) -- (g);
  \draw[<-,thick] (g) -- (m);
  \draw[<-,thick] (m) -- (h);
  \foreach \x/\t/\c in {7.1/E/avviso!40,7.5/R/white,7.9/R/white,8.3/E/avviso!40} {
    \node[draw,minimum size=3.5mm,inner sep=0pt,fill=\c] at (\x,0.55) {\t};
  }
  \node[align=center,text=torgrigio] at (7.6,1.12) {schema: quattro celle illustrative};
  \node[align=center,text=avviso] at (9,-1) {codifica il nome del servizio\\nei tipi RELAY (R) e\\RELAY\_EARLY (E)};
  \node[align=center,text=avviso] at (3,-1) {legge i tipi di cella:\\associa il client\\alla richiesta del descrittore};
  \node[align=center,text=torgrigio] at (6,-1) {inoltra i tipi di cella\\come nel caso del 2014};
\end{tikzpicture}
